\documentstyle[% 


righttag% 


,11pt% 


,amssymb% 


,russian%               remove in Eng. 


,izvru%                 Set izven% in Eng. 


% ,izvru-ks             for brief communications 


]{amsart} 


 


%       Math definitions 


 


\newcommand{\im}{\operatorname{Im}} 


%\newcommand{\ov}{\overline} 


\newcommand{\ctg}{\operatorname{ctg}} 


\newcommand{\ind}{\operatorname{ind}} 


\newcommand{\tts}[1]{} 


 


\renewcommand{\baselinestretch}{1.06}


 


%\hoffset=-15mm%                  For Eng. 


  


\setcounter{page}{62} 


\god{1998} 


\nomer{7~(434)} %                        Remove in Eng. 


%\nomer{Vol.\,42, No.\,7}%               Set in Eng. 


%\str{\pageref{begin}--\pageref{end}}%  Set in Eng. 


\udk{519.642.7} 


\author{\it Н.Я.\,Тихоненко, Е.П.\,Цымбалюк} 


% NB:   ^^ remove in Eng. 


\title{Проекционные методы решения исключительного\\ 


случая уравнений типа свертки} 


 


\date{} 


% \pagestyle{myheadings}%         Set in Eng. 


 


\pagestyle{plain} %              Remove in Eng. 


\begin{document} 


 


% \thispagestyle{plain}%          Set in Eng. 


 


\maketitle 


\label{begin} 


 


В данной работе производится обоснование проекционных методов 


приближенного решения исключительного случая уравнения типа свертки 


(УТС). В частности, рассматривается уравнение 


\begin{equation} 


\varphi(x)+\frac{1}{\sqrt{2\pi}}\int^\infty_0 


K_1(x-t)\varphi(t)dt+\frac{1}{\sqrt{2\pi}}\int^0_{-\infty} 


K_2(x-t)\varphi(t)dt=h(x),\quad x\in\Bbb R, 


\end{equation} 


где $K_1(x)$, $K_2(x)$, $h(x)$ --- известные функции, которые подчиним 


вполне определенным условиям. 


 


Приближенному решению исключительного случая УТС посвящены работы 


\cite{1}--\cite{3}, в которых устанавливаются оценки погрешности в 


пространстве $L_2$ метода приближенной факторизации решения УТС. 


Однако в 


этих работах не была установлена сходимость этого метода, также не были 


определены оценки скорости сходимости приближенных решений УТС к их 


точным решениям. 


\medskip 


 


{\bf 1${}^\circ$.} {\it Пространства.} Пусть $x\in \Bbb R$ --- вещественная 


ось; $D^+=\{z\in\Bbb C:\im z>0\}$ 


--- верхняя полуплоскость, а $D^-=\{z\in\Bbb C:\im z<0\}$ --- нижняя 


полуплоскость комплексной плоскости $\Bbb C$. 


 


Через $C^{(r)}$, где $r$ --- целое неотрицательное число, $C^{(0)}=C$, 


обозначим 


пространство $r$ раз непрерывно дифференцируемых на сомкнутой 


вещественной оси $\Bbb R$ функций $f(x)$ с нормой 


$\|f(x)\|_{C^{(r)}}=\sum\limits^r_{k=0}\max|f^{(k)}(x)|$. Через 


$C^{(r)}_{01}$, $r\ge 0$, $C^{(0)}_{01}=C_{01}$,


обозначим 


пространство функций $f(x)\in C^{(r)}$ и представимых в виде 


$f(x)=\frac{2i}{x+i}f_1(x)$. При этом 


$\|f(x)\|_{C^{(r)}_{01}}=\|f_1(x)\|_{C^{(r)}}$. Через 


$H^{(r)}_\alpha$, $r\ge 0$, $H^{(0)}_\alpha=H_\alpha$, обозначим 


пространство функций 


$f(x)\in C^{(r)}$, $r$-е производные которых удовлетворяют условию 


Г\"ельдера \cite{4} 


$$ 


|f^{(r)}(x+h)-f^{(r)}(x)|(1+|x+h|)^\alpha 


(1+|x|)^\alpha\le A h^\alpha\quad \forall x\in \Bbb R,\ \; 


\forall h>0, 


$$ 


где $A$, $\alpha$ --- постоянные; $A>0$, $0<\alpha\le 1$. Пространства 


$L$, $L^{(r)}_2$, $r\ge 0$, $L^{(0)}_2=L_2$, 


имеют обычный смысл, и в пространстве $L_2$ введены общепринятые норма 


и скалярное произведение. 


\medskip 


 


{\bf 2${}^\circ$.} {\it Аппроксимация на $\Bbb R$ рациональными функциями.} 


Пусть 


\begin{equation} 


\omega_k(x)=\biggl(\frac{x-i}{x+i}\biggr)^k,\quad 


k=0,\pm 1,\dotsc,\quad x\in\Bbb R. 


\end{equation} 


Ясно, что при $k\ge 0$ функции $\omega_k(x)$ аналитичны в $D^+$, а при 


$k\le 0$ аналитичны в 


$D^-$. Согласно \cite{5} система функций $\omega_k(x)$ ортогональна и 


полна в пространстве $L_{2\rho}$, $\rho(x)=\frac{1}{1+x^2}$. 


 


Пусть далее 


\begin{equation} 


\psi_k(x)=\biggl(\frac{x-i}{x+i}\biggr)^k\frac{2i}{x+i}, 


\quad k=0,\pm 1,\dotsc,\quad x\in \Bbb R. 


\end{equation} 


Легко видеть, что при $k\ge 0$ функции $\psi_k(x)$ аналитичны в $D^+$, а 


при $k<0$ аналитичны в $D^-$. На основании (2), (3) и работы \cite{5} 


легко 


показать, что система функций $\psi_k(x)$ ортогональна и полна в 


пространстве $L_2$. При этом $\|\psi_k(x)\|_{L_2}=2\sqrt\pi$. 


 


Обозначим через $Q_n$ оператор, который ставит в соответствие каждой 


функции $f(x)\in L_2$ отрезок ее ряда Фурье по системе функций (3),


т.\,е. 


\begin{equation} 


(Q_n f)(x)=\sum^n_{k=-n}f_k\psi_k(x), 


\end{equation} 


где 


\begin{equation} 


f_k=\frac{1}{4\pi}\int_{\Bbb R}f(x)\ov{\psi_k(x)}dx= 


\frac{1}{2\pi x}\int_{\Bbb R}f(x) 


\biggl(\frac{x+i}{x-i}\biggr)^k\frac{dx}{x-i} 


\end{equation} 


--- коэффициенты Фурье функции $f(x)\in L_2$ по системе функций (3). 


Справедливы следующие утверждения. 


 


\begin{thm} 


Справедливы оценки


$$ 


\|Q_n\|_{C_{01}\to L_2}=\|Q_n\|_{L_2\to L_2}=1,\quad 


n=1,2,\dotsc 


$$ 


\end{thm} 


 


\begin{thm} 


Если $f(x)\in L^{(r)}_2$, $r\ge 1$, или $f(x)\in C^{(r)}_{01}$, то 


$$ 


\|f(x)-(Q_n f)(x)\|_{L_2}\le d_1n^{-r}. 


$$ 


\end{thm} 


 


Здесь и далее $d_i$ --- вполне определенные постоянные, не зависящие от 


$n$. 


 


Пусть $f(x)\in C_{01}$. 


Тогда интерполяционный многочлен Лагранжа функции 


$f(x)$ по системе функций (2) согласно \cite{5} имеет вид 


\begin{equation} 


(L_n f)(x)=\sum^n_{k=-n}a_k 


\biggl(\frac{x-i}{x+i}\biggr)^k,\quad 


a_k=\frac{1}{2n+1}\sum^n_{j=-n} 


f(x_j)e^{-\frac{2\pi i}{2n+1}jk}, 


\end{equation} 


где 


\begin{equation} 


x_j=-\ctg\frac{\pi}{2n+1}j,\quad j=\ov{-n,n}. 


\end{equation} 


Легко показать, что $(L_n f)(x)\in C_{01}$ и 


$$ 


(L_n f)(x)=-\sum^n_{k=-n}\gamma_k\psi_k(x),\quad 


\gamma_k=a_n+a_{n-1}+\dotsb+a_{k+1}. 


$$ 


 


\begin{thm} 


Справедливы оценки 


$$ 


\|L_n\|_{C_{01}\to C_{01}}\le d_1\ln n,\quad 


\|L_n\|_{C_{01}\to L_2}\le 2\sqrt{2\pi},\quad n=2,3,\dotsc 


$$ 


\end{thm} 


 


\begin{thm} 


Если $f(x)\in C^{(r)}_{01}$, $r\ge 1$, то 


\begin{align*} 


\|f(x)-(L_n f)(x)\|_{C_{01}}&\le d_3n^{-r}\ln n,\\ 


\|f(x)-(L_n f)(x)\|_{L_2}&\le d_4n^{-r}. 


\end{align*} 


\end{thm} 


 


{\bf 3${}^\circ$.} {\it Разрешимость уравнения} (1). Пусть 


$k_1(x),k_2(x)\in L$, а $h(x)\in L_2$. Тогда 


согласно \cite{6} уравнение (1) эквивалентно следующей задаче Римана 


\begin{equation} 


(K\Phi)(x)\equiv[1+K_1(x)]\Phi^+(x)- 


[1+K_2(x)]\Phi^-(x)=H(x),\quad x\in \Bbb R, 


\end{equation} 


где $K_1(x)$, $K_2(x)$, $H(x)$ --- преобразования Фурье \cite{7} 


соответственно функций 


$k_1(x)$, $k_2(x)$, $h(x)$. При этом решение уравнения (1) и задачи 


Римана (8) связаны следующим образом 


\begin{equation} 


\varphi(x)=\frac{1}{\sqrt{2\pi}}\int_{\Bbb R}[\Phi^+(t)-\Phi^-(t)] 


e^{-ixt}dt. 


\end{equation} 


 


Пусть $K_1(x),K_2(x)\in H^{(r+1)}_\alpha$, $0<\alpha\le 1$, 


$r\ge r_0$, где 


число $r_0$ определим ниже, и функции $1+K_1(x)$, $1+K_2(x)$ имеют 


нули соответственно в точках 


$\alpha_0,\alpha_1,\dotsc,\alpha_p$; 


$\beta_0,\beta_1,\dotsc,\beta_q$ на вещественной оси $\Bbb R$, 


$\alpha_0=\infty$, $\beta_0=\infty$, 


соответственно целых порядков 


$\nu_0,\nu_1,\dotsc,\nu_p$; $\mu_0,\mu_1,\dotsc,\mu_q$. Тогда согласно 


\cite{8} справедливы представления 


\begin{equation} 


1+K_1(x)=A(x)\rho_-(x),\quad 


1+K_2(x)=B(x)\rho_+(x), 


\end{equation} 


где 


\begin{align} 


\rho_-(x)&=(x-i)^{-\nu_0} 


\biggl(\frac{x-\alpha_1}{x-i}\biggr)^{\nu_1}\dotsb 


\biggl(\frac{x-\alpha_p}{x-i}\biggr)^{\nu_p},\\ 


\rho_+(x)&=(x-i)^{-\mu_0} 


\biggl(\frac{x-\beta_1}{x-i}\biggr)^{\mu_1}\dotsb 


\biggl(\frac{x-\beta_q}{x-i}\biggr)^{\mu_q},


\end{align} 


а функции $A(x),B(x)\in H^{(r+1)}_\alpha$, $0<\alpha\le 1$, 


$r\ge r_0$; $A(x)\ne 0$, $B(x)\ne 0$, 


где 


\begin{equation} 


r_0=\max\{\nu_0,\nu_1,\dotsc,\nu_p;\, 


\mu_0,\mu_1,\dotsc,\mu_q\}. 


\end{equation} 


Тогда задачу Римана (8) можно переписать в виде 


\begin{equation} 


K\Phi\equiv(A\rho_-P+B\rho_+Q)\Phi=H, 


\end{equation} 


где $P=\frac{1}{2}(I+S)$, $Q=\frac{1}{2}(Q+S)$ --- проекторы Рисса; $I$ 


--- единичный оператор, а $S$ --- сингулярный интеграл Коши. 


 


Заметим, что случай $\nu_0\ne 0$, $\mu_0\ne 0$ соответствует случаю, 


когда уравнение (1) является уравнением первого рода. 


 


Пусть $\varkappa=\ind[B(x)/A(x)]$. Тогда при данных предположениях 


относительно функций $K_1(x)$ 


и $K_2(x)$ при $\varkappa\ge 0$ задача Римана (8), а значит, и уравнение 


(1), безусловно 


разрешимы, если $H(x)\in L^{(r)}_2$, $r\ge 0$, и решение задачи 


принадлежит $L_2$. Справедлива 


\medskip


 


{\bf Лемма} \cite{9}. {\it


Пусть функция $C(x)\in H^{(r+1)}_\alpha$, $r\ge 0$, 


$\frac{3}{4}<\alpha\le 1$. Тогда коммутатор 


$T:L_2\to C^{(r)}_{01}$, $r\ge 0$, компактен, где


}


$$ 


(T\varphi)(x)=\frac{1}{\pi i}\int_{\Bbb R} 


\frac{C(t)-C(x)}{t-x}dt. 


$$ 


 


В дальнейшем приближенные решения уравнения (1) будем строить на основе 


приближенного решения задачи (8). 


\medskip 


 


{\bf 4${}^\circ$.} {\it Метод редукции.} Пусть $\frak M$ --- алгебра 


элементов 


$C(x)\in H^{(r+1)}_\alpha$, $r\ge r_0$, $\frac{3}{4}<\alpha<1$, таких, 


что $C(x)\ne 0$ на 


$\Bbb R$, $\ind C(x)=0$. Тогда согласно \cite{6}, \cite{8} справедливы 


представления 


$$ 


C(x)=C_+(x)C_-(x),\quad C(x)=d_+(x)d_-(x), 


$$ 


где $C^{\pm 1}_\pm(x)\ne 0$, $d^{\pm 1}_\pm(x)\ne 0$ на $\Bbb R$; 


$C^{\pm 1}_\pm(x)$, $d^{\pm 1}_\pm(x)$ аналитичны соответственно в 


$D^\pm$, $C^{\pm 1}_\pm(x),\linebreak d^{\pm 1}_\pm(x)\in H^{(r+1)}_\alpha$, 


$r\ge r_0$, $\frac{3}{4}<\alpha< 1$, т.\,е. 


$C^{\pm 1}_\pm(x),d^{\pm 1}_\pm(x)\in\frak M$. 


 


Приближенное решение задачи (8) ищем в виде 


\begin{equation} 


\Phi_n(x)=\sum^n_{k=-n}\varphi_k\psi_k(x), 


\end{equation} 


где функции $\psi_k(x)$ имеют вид (3). Ясно, что 


$\Phi_n(x)=\Phi^+_n(x)-\Phi^-_n(x)$, где 


\begin{equation} 


\Phi^+_n(x)=\sum^n_{k=0}\varphi_k\psi_k(x),\quad 


\Phi^-_n(x)=-\sum^{-1}_{k=-n}\varphi_k\psi_k(x). 


\end{equation} 


Неизвестные постоянные $\varphi_k$ определяем из системы уравнений 


\begin{equation} 


\sum^n_{k=0}A_{jk}\varphi_k-\sum^{-1}_{k=-n} 


B_{jk}\varphi_k=H_j,\quad j=\ov{-n,n}, 


\end{equation} 


где $A_{jk}$, $B_{jk}$, $H_j$ --- коэффициенты Фурье (5) соответственно 


функций $[1+K_1(x)]\psi_k(x)$,\linebreak $[1+K_2(x)]\psi_k(x)$, $H(x)$


по системе функций (3). 


 


\begin{thm} 


Пусть $k_1(x),k_2(x)\in L;$ $K_1(x),K_2(x)\in H^{(r+1)}_\alpha$, 


$\frac{3}{4}<\alpha<1$, $r\ge r_0$, 


где число $r_0$ определено соотношением $(13);$ 


справедливы представления $(10)$, где $A(x)\ne 0$, $B(x)\ne 0$, 


$\ind A(x)=\ind B(x)=0;$ $H(x)\in L^{(r)}_2$, 


$r\ge r_0$. Тогда система 


уравнений $(17)$ однозначно разрешима при достаточно больших $n$, и 


приближенные решения $\Phi_n(x)$ задачи $(8)$ сходятся в пространстве 


$L_2$ к ее точному решению $\Phi(x)$ со скоростью 


\begin{equation} 


\|\Phi(x)-\Phi_n(x)\|_{L_2}\le d_5n^{-r}. 


\end{equation} 


\end{thm} 


 


\begin{pf} 


Заметим сначала, что в условиях теоремы задача (8) имеет единственное 


решение, принадлежащее пространству $L_2$. Далее будем следовать 


результатам главы XI работы \cite{8}. Положим $X=Y=L_2$. Пусть $P_n$ --- 


оператор, который ставит в соответствие каждому элементу пространства 


$X$ агрегат вида (15), (16), а оператор $Q_n$ 


--- оператор (4). Ясно, что 


эти операторы обладают свойствами 


$\im P_n\subset X$, $\im Q_n\subset Y$, $P^2_n=P_n$, $Q^2_n=Q_n$, 


$\dim\im P_n=\dim\im Q_n=2n+1$. 


Тогда систему (17) можно записать в виде 


\begin{equation} 


(Q_n K P_n\Phi)(x)=(Q_n H)(x). 


\end{equation} 


Ясно, что в условиях теоремы оператор $K$ непрерывно обратим. Покажем 


теперь, что система (19), а значит, и система (17) разрешимы при 


достаточно больших $n$, а приближенные решения задачи (8) сходятся к ее 


точным решениям по норме пространства $X$; для этого достаточно 


показать, что к оператору $K$ применим \cite{8} проекционный процесс по 


системе проекторов $P_n$ и $Q_n$. Рассмотрим уравнение 


$$ 


B\Phi\equiv(P\rho_-+Q\rho_+)\Phi=0, 


$$ 


где $P$, $Q$ --- проекторы Рисса, а $\rho_\pm(x)$ определяется по 


формулам (11), (12). Союзное ему уравнение является задачей Римана 


$$ 


\rho_-(x)R^+(x)=\rho_+(x)R^-(x),\quad x\in\Bbb R, 


$$ 


которая \cite{7} имеет лишь тривиальное решение. Поэтому 


$\dim\ker B=0$. Рассмотрим 


теперь пространства $Z=L_2$, $Z_0=C^{(r)}_{01}$, $r\ge r_0$, и построим 


пространство $\ov X=\{\Phi(x):\Phi(x)\in X,\ (B\Phi)(x)\in Z\}$, которое 


становится банаховым, если в нем ввести норму 


$\|\Phi(x)\|_X=\|(B\Phi)(x)\|_X$. Легко видеть, что 


оператор $P$ ограничен \cite{5} в пространстве $Z$, и оператор умножения 


элементов пространства $Z_0$ на элементы алгебры $\frak M$ 


также является 


ограниченным, а потому следует включение $Z\subset\ov X$. Следовательно, 


сужение $\ov B$ 


оператора $B$ на пространство $\ov X$ взаимнооднозначно и непрерывно 


отображает пространство $\ov X$ на пространство $Z$, т.\,е. 


$\ov B:\ov X\to Z$ непрерывно 


обратим\,\footnote{$^)$Здесь и далее ниже черта над оператором означает 


его сужение над соответствующим пространством.}{$^)$}. 


 


Так как $A(x)\ne 0$, $B(x)\ne 0$ на $\Bbb R$, 


$\ind A(x)=\ind B(x)=0$, $A(x),B(x)\in H^{(r+1)}_\alpha$, 


$r\ge r_0$, $\frac{3}{4}<\alpha< 1$, то 


справедливы \cite{6}, \cite{8} факторизации 


\begin{equation} 


A(x)=A_+(x)A_-(x),\quad B(x)=B_+(x)B_-(x), 


\end{equation} 


где $A^{\pm 1}_\pm(x)\ne 0$, $B^{\pm 1}_\pm(x)\ne 0$ на $\Bbb R$; 


$A^{\pm 1}_\pm(x)$, $B^{\pm 1}_\pm (x)$ 


аналитичны в $D^\pm$, 


$A^{\pm 1}_\pm(x),B^{\pm 1}_\pm(x)\in H^{(r+1)}_\alpha$, 


$r\ge r_0$, $\frac{3}{4}<\alpha< 1$. Пусть $E:\ov X\to X$ 


--- оператор вложения. Ясно, что он ограничен. Тогда для сужения 


$\ov K$ 


оператора $K$ справедливо \cite{8} представление 


$$ 


\ov K=K_1+T_1, 


$$ 


где 


$$ 


K_1=(A_+P+B_-Q)(PA_-+QB_+)\ov B,\quad 


T_1=(A_+Q A_-\rho_-P+B_-P B_+\rho_+Q)E. 


$$ 


На основании \cite{6} операторы $A_+P+B_-P$, $PA_-+QB_+$ 


обратимы, т.\,к. представляют 


собой сингулярные операторы с нулевыми индексами и отличными от нуля 


символами, а поэтому оператор $K_1:\ov X\to Z$ 


обратим как произведение обратимых 


операторов. На основании (20) и леммы оператор 


$A_+QA_-\rho_-P+B_-PB_+\rho_+Q:X\to Z_0$ компактен. Поэтому 


оператор $T_1:\ov X\to Z_0$ компактен как произведение компактного и 


ограниченного 


операторов. Пусть $X^\pm(x)\in\frak M$. Тогда на основании работы 


\cite{5} в силу (2) справедливы разложения 


\begin{equation} 


X^+(x)=\sum^\infty_{k=0}\eta^+_k\biggl(\frac{x-i}{x+i}\biggr)^k, 


\quad 


X^-(x)=\sum^0_{k=-\infty}\eta^-_k\biggl(\frac{x-i}{x+i}\biggr)^k. 


\end{equation} 


Тогда на основании представлений (21), (15) следует цепочка 


равенств 


\begin{align} 


&P_n[PX^-+QX^+]P_n\Phi=P_n[PX^-+QX^+]\Phi_n= 


P_n[PX^-\Phi_n+QX^+\Phi_n]=\notag \\ 


%


&\quad =P_n\biggl[P\sum^0_{k=-\infty}\eta^-_k 


\biggl(\frac{x-i}{x+i}\biggr)^k\sum^n_{k=-n}\varphi_k\psi_k(x)+ 


Q\sum^\infty_{k=0}\eta^+_k\biggl(\frac{x-i}{x+i}\biggr)^k 


\sum^n_{k=-n}\varphi_k\psi_k(x)\biggr]=\notag \\ 


%


&\quad =P_n\biggl[P\sum^n_{k=-\infty}A_k\psi_k(x)+


Q\sum^\infty_{k=-n}B_k\psi_k(x)\biggr]= 


P_n\biggl[\sum^n_{k=0}A_k\psi_k(x)+


\sum^{-1}_{k=-n}B_k\psi_k(x)\biggr]=\notag \\


%


&\quad=\sum^n_{k=0}A_k\psi_k(x)+\sum^{-1}_{k=-n} 


B_k\psi_k(x)=\sum^n_{k=-n}C_k\psi_k(x).


\end{align} 


Из этой цепочки вытекает равенство 


$P_n[PX^-+QX^+]P_n=[PX^-+QX^+]P_n$, из которого в силу 


конечномерности операторов $P_n$ следует включение 


$\im P_n\subset B(\im P_n)$. Легко видеть, что 


$\im P_n=\im Q_n\subset Z$. Тогда на основании теоремы о непрерывном 


графике сужение $\ov P_n$ оператора $P_n$ на пространство $\ov X$ 


является непрерывным проектором в пространстве $X$. При этом в силу 


определения пространств $Z$ и $Z_0$ на 


основании теоремы 1 оператор $Q_n:Z_0\to Z$ является ограниченным. В 


качестве 


пространства $Z_1$ положим $L^{(r)}_2$, $r\ge r_0$. Тогда очевидно 


включение $Z_0\subset Z_1\subset B(Q_n)$. При 


этом на основании теоремы 2\; $(Q_n\Phi)(x)\to\Phi(x)$ по норме 


пространства $X$ для любого элемента $\Phi(x)\subset Z_1$. Рассмотрим 


теперь оператор $C=X^+P+X^-Q$, 


где $X^\pm(x)\ne 0$ на $\Bbb R$, $X^\pm(x)$ аналитичны 


в $D^\pm$, $X^\pm(x)\in H^{(r+1)}_\alpha$, $r\ge r_0$, 


$\frac{3}{4}<\alpha< 1$. По аналогии с 


цепочкой (22) легко проверяется 


равенство $Q_n CQ_n=CQ_n$. Кроме того, легко проверить, что 


$C(Z_1)\subset Z_1$, $C(Z_0)\subset Z_0$. 


Тогда на 


основании теоремы 1.4 главы XI работы \cite{8} следует, что пространство 


$Z_1$ принадлежит многообразию сходимости оператора $K_1$ по системе 


проекторов $\ov P_n$ и $Q_n$. Это означает, что при достаточно больших 


$n$ существуют операторы $(Q_n K_1\ov P_n)^{-1}$, и последовательность 


$x_n=(Q_n K_1\ov P_n)^{-1}Q_n y$ для любого $y\in Y$ сходится 


но норме пространства $X$ к некоторому элементу из $X$, т.\,е. к 


оператору $K_1$ примен\'им проекционный процесс по системе проекторов 


$\ov P_n$ и $Q_n$. Так как оператор $T_1$ компактен и оператор $\ov K$ 


обратим, то на 


основании теоремы 1.1 главы XI работы \cite{8} следует, что к оператору 


$\ov K=K_1+T_1$ также примен\'им проекционный процесс по системе 


проекторов $\ov P_n$ и $Q_n$. 


Это означает, что система (17) разрешима при достаточно больших $n$. 


Легко видеть, что 


$$ 


y_0(x)=(A^{-1}_+P+B^{-1}_-Q)[H-T_1\Phi](x)\in L^{(r)}_2,\quad r\ge r_0. 


$$ 


Тогда на основании теоремы 2.3 главы XI работы \cite{8} из теоремы 2 


следует оценка (18). 


\end{pf} 


 


Так как система (17) разрешима при достаточно больших $n$, то при таких 


$n$ постоянные $\varphi_k$ определяются однозначно. Тогда приближенные 


решения уравнения (1) строятся по формуле 


\begin{equation} 


\varphi_n(x)=\frac{1}{\sqrt{2\pi}}\int_{\Bbb R} 


[\Phi^+_n(t)-\Phi^-_n(t)]e^{-ixt}dt, 


\end{equation} 


где $\Phi^\pm_n(x)$ --- приближенные решения задачи (8). Согласно 


\cite{7}, \cite{6} 


приближенные решения уравнения (1) в явном виде замкнутся следующим 


образом 


\begin{equation} 


\varphi_n(x)= 


\begin{cases} 


\displaystyle 2\sqrt{2\pi}e^{-x}\sum^n_{k=0}\varphi_k 


\sum^k_{j=0}C^j_k\frac{2^j}{j\,!}(-1)^j x^j, & x>0;\\[2mm] 


\displaystyle 2\sqrt{2\pi}e^x\sum^{-1}_{k=-n}\varphi_k 


\sum^{-1}_{j=0}C^j_{-k}\frac{2^j}{j\,!}x^j, &x<0. 


\end{cases} 


\end{equation} 


Тогда на основании представлений (9), (23), равенства Парсеваля для 


интегралов Фурье \cite{7} и теоремы 5 справедлива 


 


\begin{thm} 


Пусть функции $k_1(x),k_2(x)\in L;$ функции 


$K_1(x),K_2(x)\in H^{(r+1)}_\alpha$, $\frac{3}{4}<\alpha< 1$, 


$r\ge r_0$, где 


число $r_0$ определено 


соотношением $(13);$ справедливы представления $(10)$, где 


$A(x)\ne 0$, $B(x)\ne 0$, $\ind A(x)=\ind B(x)=0$, 


$H(x)\in L^{(r)}_2$, $r\ge r_0$. Тогда система уравнений $(17)$ 


однозначно разрешима при 


достаточно больших $n$, а приближенные решения уравнения $(1)$ строятся 


по формулам $(24)$, где $\varphi_k$ --- решения системы $(17)$, и 


сходятся в пространстве $L_2$ к его точному решению со скоростью 


\begin{equation} 


\|\varphi(x)-\varphi_n(x)\|_{L_2}\le d_5n^{-r}. 


\end{equation} 


\end{thm} 


 


{\bf 5${}^\circ$.} {\it Метод коллокаций.} Приближенное решение задачи (8) 


ищем в виде (15), (16), а неизвестные постоянные $\varphi_k$ 


определяем из системы уравнений 


\begin{equation} 


\sum^n_{k=0}[1+K_1(x_j)]\psi_k(x_j)\varphi_k+ 


\sum^{-1}_{k=-n}[1+K_2(x_j)]\psi_k(x_j)\varphi_k=H(x_j), 


\quad j=\ov{-n,n}, 


\end{equation} 


где $x_j$ --- узлы (7). Ясно, что систему (26) можно записать в виде 


$$ 


(L_n K P_n\Phi)(x)=(L_n H)(x), 


$$ 


где оператор $P_n$ имеет прежний смысл, а оператор $L_n$ --- оператор 


(6). 


Здесь $X=L_2$, $Y=C_{01}$, $Z=L_2$, $Z_0=C^{(r)}_{01}$, 


$Z_1=L^{(r)}_2$, $r\ge r_0$. Тогда оператор $L_n$ обладает свойством 


$L^r_n=L_n$. Справедлива 


 


\begin{thm} 


Пусть функции $k_1(x),k_2(x)\in L$, функции 


$K_1(x),K_2(x)\in H^{(r+1)}_\alpha$, $\frac{3}{4}<\alpha< 1$, 


$r\ge r_0$, где число $r_0$ определено 


соотношением $(13);$ справедливы представления $(10)$, где 


$A(x)\ne 0$, $B(x)\ne 0$, $\ind A(x)=\ind B(x)=0;$ 


$H(x)\in C^{(r)}_{01}$, $r\ge r_0$. Тогда система $(26)$ разрешима 


при достаточно больших $n$, а 


приближенные решения уравнения $(1)$ строятся по формулам $(24)$, где 


$\varphi_k$ --- решения системы $(26)$, и сходятся в пространстве $L_2$ 


к его точному решению со скоростью $(25)$. 


\end{thm} 


 


{\bf 6${}^\circ$.} {\it Уравнение Винера--Хопфа.} Если $k(x)\in L$, а 


$h(x)\in L_2$, $h(x)\equiv 0$ при $x<0$, то согласно \cite{6} уравнение 


Винера--Хопфа 


\begin{equation} 


\varphi(x)+\frac{1}{\sqrt{2\pi}}\int^\infty_0 


k(x-t)\varphi(t)dt=h(x),\quad x>0 


\end{equation} 


эквивалентно задаче Римана 


\begin{equation} 


[1+K(x)]\Phi^+(x)-\Phi^-(x)=H(x),\quad x\in \Bbb R, 


\end{equation} 


где $K(x)$, $H(x)$ --- преобразования Фурье соответственно функций 


$k(x)$ и $h(x)$. 


При этом решение уравнения (27) выражается через решение задачи (28) 


следующим образом 


$$ 


\varphi(x)=\frac{1}{\sqrt{2\pi}}\int_{\Bbb R} 


\Phi^+(t)e^{-ixt}dt,\quad x>0. 


$$ 


 


Пусть $K(x)\in H^{(r+1)}_\alpha$, $0<\alpha<1$, $r\ge r_0$, где число 


$r_0$ определим ниже, и функция $1+K(x)$ имеет нули 


в точках $\alpha_0,\alpha_1,\dotsc,\alpha_p$ на вещественной оси 


$\Bbb R$, $\alpha_0=\infty$, 


соответственно кратностей $\nu_0,\nu_1,\dotsc,\nu_p$. 


Тогда справедливо представление $1+K(x)=A(x)\rho_-(x)$, где $\rho_-(x)$ 


имеет вид (12), $A(x)\in H^{(r+1)}_\alpha$, $0<\alpha<1$, $r\ge r_0$; 


$A(x)\ne 0$, где $r_0=\max\{\nu_0,\nu_1,\dotsc,\nu_p\}$. Приближенные 


решения задачи Римана (28) ищем в виде (15), 


(16), а неизвестные постоянные $\varphi_k$ в случае метода редукции 


определяем из системы уравнений 


\begin{equation} 


\sum^n_{k=0}A_{jk}\varphi_k+\gamma_0\varphi_j=H_j,\quad 


j=\ov{-n,n}, 


\end{equation} 


где $\gamma_0=0$ при $j\ge 0$, $\gamma_0=1$ при $j\le 0$, а $A_{kj}$, 


$H_j$ --- коэффициенты Фурье 


соответственно функций $[1+K(x)]\psi_k(x)$, $H(x)$. В случае метода 


коллокаций неизвестные постоянные $\varphi_k$ определяем из системы 


уравнений 


\begin{equation} 


[1+K(x_j)]\sum^n_{k=0}\psi_k(x_j)\varphi_k+ 


\sum^{-1}_{k=-n}\psi_k(x_j)\varphi_k=H(x_j),\quad 


j=\ov{-n,n}, 


\end{equation} 


где $x_j$ --- узлы (7). Приближенные решения $\varphi_n(x)$ уравнения 


(27) строятся по формуле 


$$ 


\varphi_n(x)= 


\begin{cases} 


\displaystyle 2\sqrt{2\pi}e^{-x}\sum^n_{k=0}\varphi_k 


\sum^k_{j=0}C^j_k\frac{2^j}{j\,!}(-1)^j x^j, & x>0;\\[2mm] 


0, & x<0. 


\end{cases} 


$$ 


 


По аналогии с теоремами 5--7 устанавливается разрешимость при достаточно 


больших $n$ систем (29), (30) и определяются оценки скорости сходимости 


приближенные решений уравнения (27) к его точному решению. 


 


Заметим также, что случай $\nu_0\ne 0$ соответствует случаю уравнения 


Винера--Хопфа первого рода. 


 


По изложенной выше схеме строятся приближенные решения исключительного 


случая парного уравнения 


$$ 


\begin{cases} 


\displaystyle \varphi(x)+\frac{1}{\sqrt{2\pi}}\int_{\Bbb R} 


k_1(x-t)\varphi(t)dt=h(x), &x>0;\\[4mm] 


\displaystyle \varphi(x)+\frac{1}{\sqrt{2\pi}}\int_{\Bbb R} 


k_2(x-t)\varphi(t)dt=h(x), &x<0. 


\end{cases} 


$$ 


 


\begin{thebibliography}{9} 


 


\bibitem{1} 


Черський Ю.Й. {\em Про наближене розв"язання рIвняння Вiнера--Хопфа 


першого роду} // ДАН УРСР. -- 1966. -- \No\,8. -- С.\,992--995. 


\bibitem{2} 


Черский Ю.И. {\em Приближенное решение уравнения Винера--Хопфа в одном 


исключительном случае} // Дифференц. уравнения. -- 1966. -- Т.\,2. -- 


\No\,8. -- С.\,1093--1100. 


\bibitem{3} 


Тихоненко М.Я. {\em До наближеного розв"язку iнтегральних та 


iнтегро-диференцiальних рiвнянь з рiзницевими ядрами в узагальнених 


функцiях} // ДАН УРСР. -- 1972. -- \No\,6. -- С.\,536--539. 


\bibitem{4} 


Самко С.Г., Килбас А.А., Маричев О.И. {\em Интегралы и производные 


дробного порядка и некоторые их приложения}. -- Минск, 1987. -- 687~с. 


\bibitem{5} 


Иванов В.В. {\em Теория приближенных методов и ее применение к 


численному решению сингулярных интегральных уравнений}. -- Киев: Наук. 


думка, 1968. -- 287~с. 


\bibitem{6} 


Гахов Ф.Д., Черский Ю.И. {\em Уравнения типа свертки}. -- М.: Наука, 


1978. -- 296~с. 


\bibitem{7} 


Князев П.Н. {\em Интегральные преобразования}. Учеб. пособие. -- Минск: 


Вышэйш. школа, 1969. -- 197~с. 


\bibitem{8} 


Пресдорф З. {\em Некоторые классы сингулярных уравнений}. -- М.: Мир, 


1979. -- 493~с. 


\bibitem{9} 


Свяжина Н.Н., Тихоненко Н.Я. {\em Аппроксимация функций на вещественной 


оси и проекционные методы решения исключительного случая сингулярных 


интегральных уравнений}. -- Одес.\linebreak ун-т. -- Одесса, 1995. -- 69~с. -- 


Деп. в ГНТБ Украины 16.02.95, \No\,383.-Ук.\,95. 


 


 


\end{thebibliography} 


 


\vskip 0.5cm 


 


\post{\it Одесский государственный}{\it Поступила} 


\post{\it университет $($Украина$)$}{10.07.1995} 


\smallskip 


\post{\it Украинская государственная}{} 


\post{\it академия связи}{} 


 


\label{end} 


\end{document} 


